\chapter{An Element's Identity Card}

\begin{kaichang}
You probably have a pencil handy. Its barrel may say ``HB'' or ``2B.'' In Chinese it is called a ``lead pen,'' and English also calls it a lead pencil. Here is a slightly disappointing fact, though: pencil cores contain \textbf{not the slightest bit of lead}. They are pressed from graphite and clay, and graphite is pure carbon.

Lead in name, carbon in fact. What gets the final say? What actually determines what an element is? Consider an older obsession too: for hundreds of years, countless alchemists tried to turn dull gray lead into gleaming gold. They stoked furnaces, added potions, and chanted spells, one after another, and all failed. Were they unskilled, or was the enterprise fundamentally impossible?

Write down your guesses first. By the end of this chapter, we will issue every element an identity card, then see straight through both questions.
\end{kaichang}

\section{The Cards Chemical Reactions Cannot Change}

Begin with observation. Nails rust, copper pots turn green, and sugar heated too long turns black: matter changes its appearance in countless ways. Seventeenth- and eighteenth-century chemists spent more than two hundred years burning, boiling, dissolving, and electrifying substances and discovered something strange. Some things could never be turned into something else, however they were treated.

Heat iron red-hot, hammer it into sheets, grind it into powder, or dissolve it in acid, and iron remains iron. React it separately with oxygen, sulfur, and chlorine, and iron can still be found in the products. Draw gold into wire much thinner than a hair, dissolve it in aqua regia, and find a way to precipitate it again: it is still gold. Lavoisier's generation used this list of things that survived every treatment to define elements: the basic components that chemical reactions can neither decompose nor transform into one another.

In other words, a chemical reaction shuffles a deck of cards. Cards change combinations and partners, but every card remains. Burn charcoal, and its carbon has not disappeared; it has simply changed partners and moved into the air. Chapter 26 will account for this in detail.

Daily life echoes this list everywhere. A gold ring is still gold after ten years, though a layer may have worn away. Iodine in iodized salt remains iodine when dissolved in soup and eaten. Nitrogen in a breath visits your lungs and leaves unchanged: your body cannot break it apart. Without this list, ancient people mistook turning stone into gold for a mere technical problem. The question therefore moves down a level: beneath the face of the card, what locks an element's identity in place?

\section{Proton Number: The Identity Number}

Chapter 6 introduced the atom's three main characters: positively charged protons and uncharged neutrons crowded into a nucleus tens of thousands of times smaller than the atom, with negatively charged electrons outside. Which determines elemental identity?

Eliminate them one at a time. Electrons are the first suspects and the first ruled out. Sodium atoms in a sodium lamp and sodium ions in salt have different electron counts, 11 and 10 respectively, yet chemists call both sodium. Electron number clearly does not decide. Neutrons? Shortly we will see that different neutron counts can also belong to one element. That leaves protons: they supply the identity number.

\textbf{An element is the class of all atoms with the same number of protons in their nuclei.} Proton number has its own name: \textbf{atomic number}. Hydrogen's is 1, carbon's 6, oxygen's 8, iron's 26, and gold's 79. Every nucleus with 6 protons is carbon, regardless of its surrounding electrons, its neutron count, or whether it is inside your body or embedded in a diamond. Lead has 82 protons and gold has 79. Three protons separate them: the entire reason lead cannot become gold.

This seems obvious today, but atomic number originally meant only an element's position in a table. Nobody knew what physical feature inside an atom it represented. In 1913, the young British physicist Moseley bombarded elements with X-rays and found that the frequency of their emitted X-rays changed strictly with integer steps. The integer was precisely the atomic number and could be interpreted directly as the number of positive nuclear charges. Atomic number gained a physical meaning: nuclear proton count. Incidentally, this vindicated Mendeleev. His insistence on placing tellurium before iodine according to their personalities, opposite their atomic-weight order, had attracted ridicule. Measured atomic numbers showed he was exactly right (Chapter 10 returns to the story). Moseley died fighting at Gallipoli in 1915, aged only 27. Historians of science often mourn him as one of the most poignant examples of a scientist who never received a Nobel Prize.

Why should proton number decide? Follow the causal chain. Nuclear protons carry positive charge, whose amount determines how many electrons a neutral atom can hold. Chemical behavior---preferred reaction partners and ways of reacting---depends almost entirely on electron number and arrangement (Chapters 9 and 11 unpack this). Proton number is therefore furthest upstream: it determines the electron household, which determines personality. Neutrons hide uncharged inside the nucleus and have no say over electrons; electrons themselves can come and go. Through the entire performance, only protons never change actors.

\section{A Few More Neutrons: Isotopes}

The identity number fixes proton count, but nuclei sharing that number can have different weights because neutron count can vary.

Take carbon. The sum of protons and neutrons is called the \textbf{mass number}. Most carbon nuclei have 6 protons and 6 neutrons, mass number 12, written carbon-12. About 1 in every 100 has one extra neutron: carbon-13. Extremely rare carbon-14 has 6 protons and 8 neutrons. All have 6 protons, so all are carbon, with nearly identical chemical behavior but different nuclear weights. Atoms of the same element with \textbf{equal proton counts and different neutron counts} are \textbf{isotopes} of one another. The term means ``same place'' in Greek, because they share a periodic-table box (Chapter 10 unfolds the full table).

Hydrogen shows the most dramatic isotope differences. Ordinary hydrogen has only 1 proton in its nucleus; deuterium adds 1 neutron, doubling the weight; tritium adds 2, tripling it. How significant is doubling? Carbon-13 is only about 8\% heavier than carbon-12, but deuterium is a full 100\% heavier than ordinary hydrogen. Water containing deuterium, heavy water, is therefore about 10\% denser than ordinary water and freezes nearly 4 degrees higher: ordinary water freezes at $0\,^{\circ}\mathrm{C}$, heavy water at about $3.8\,^{\circ}\mathrm{C}$. Different weights also slow reactions, the isotope effect. Enzymes in living things show a slight preference among lighter and heavier isotopes. Scientists measure that preference to infer ancient diets and whether past climates were cold or warm.

Most isotopes are stable: carbon-12, carbon-13, oxygen-16, and deuterium can remain unchanged for billions of years. A minority have unfavorable proton--neutron proportions and spontaneously change at some moment; carbon-14 and tritium are on that list. Stable isotopes already permeate everyday life. Heavy water works in some nuclear reactors (Chapter 8 reveals its job), and oxygen-18-to-oxygen-16 ratios in deep glacier ice cores vary with ancient temperatures, providing one of climatologists' most trusted thermometers. How unstable isotopes change and what they become is a question we will hold for Chapter 8.

\section{A Few More Electrons: An Ion Remains the Same Element}

Now consider electrons. In an electrically neutral atom, electron count exactly equals proton count. But electrons occupy the atom's outside and are the easiest of its three particles to move in or out. When sodium meets chlorine, one sodium electron transfers to chlorine. Sodium becomes positive \chem{Na^+}, chlorine negative \chem{Cl^-}. A charged atom or group of atoms is an \textbf{ion}.

Watch what changes and what stays. Electron count and total charge change, and chemical personality transforms: silvery sodium metal, violently reactive with water, disappears, as does poisonous yellow-green chlorine gas. Yet not one nuclear particle on either side moves: 11 protons remain 11 and 17 remain 17. Thus \chem{Na^+} is still sodium and \chem{Cl^-} still chlorine, the same elements in different charge states. That is why sodium chloride, table salt, can be eaten safely. Chlorine in salt and chlorine in poisonous gas are the same element, but the same element never means the same substance (Chapter 10 also emphasized this through salt and disinfectant).

You can see this on supermarket shelves. Sports-drink labels list sodium and potassium, but no manufacturer would dare put the metals in a bottle: that would be lethal. They contain sodium and potassium ions, the stable forms after those elements surrender electrons. Calcium supplements likewise contain calcium ions in calcium carbonate, not grayish white pieces of calcium metal. Sellers speak of elements; bottles contain ions. Consumers may be confused, but chemists distinguish them clearly.

A small table brings this together:

\begin{table}[htbp]
\centering
\begin{tabular}{lll}
\toprule
What changes & Elemental identity & Result \\
\midrule
Electron count & Unchanged & Ion; very different chemistry \\
Neutron count & Unchanged & Isotope; almost the same chemistry \\
Proton count & Changed & Another element; impossible chemically \\
\bottomrule
\end{tabular}
\end{table}

\section{Writing the Identity Card in Symbols}

Now the symbols can enter. One or two letters represent an element: \chem{C} for carbon, \chem{O} for oxygen. Some derive from Latin names: iron is \chem{Fe}, from ferrum, rather than a match to its Chinese name. A full identity card includes number and weight: atomic number, the proton count, at lower left and mass number at upper left. For example,

\[{}^{12}_{6}\chem{C} \qquad {}^{14}_{6}\chem{C} \qquad {}^{2}_{1}\chem{H}\]

Reading is straightforward: ${}^{14}_{6}\chem{C}$ means element 6, mass number 14---carbon-14. Neutron number need not be written: subtract lower left from upper left, $14-6=8$ neutrons. For ions, put net charge at upper right: \chem{Na^+}, \chem{Cl^-}, \chem{Ca^{2+}}. The number counts electrons lost or gained.

\section{Weighing Atoms}

\begin{zhengju}
We just said most carbon is carbon-12 and about 1\% is carbon-13. How do scientists know? They cannot inspect atoms one by one. The answer is an instrument designed to weigh atoms, whose origins involve a real mystery.

By the late nineteenth century, chemists had measured relative atomic masses and noticed an awkward fact: many were not integers. Neon, for example, repeatedly came out around 20.2. If every element consisted of identical atoms as Dalton supposed, where did the 0.2 come from? Some suspected measurement error, others impurity in natural elements, and some doubted atomic theory itself.

In 1913, Thomson sent neon ions through electric and magnetic fields. Under the old picture, equal masses should leave just one parabola on a photographic plate. Two appeared instead: one corresponding to mass 20 and a fainter one to mass 22. One element, two weights. His student Aston saw the instrument's promise and in 1919 built the first truly precise \textbf{mass spectrometer}. Its idea had three steps: turn atoms into charged ions, accelerate them electrically, then pass them through a magnetic field. The field bends moving charged particles, bending lighter ones more easily than heavier ones. Different masses reach different positions, like balls of different weights rolling through a crosswind. Aston examined many elements and confirmed neon as roughly 90\% neon-20 and 10\% neon-22, giving a weighted average of exactly 20.2. The noninteger-mass mystery was solved: the hidden assumption of one kind of atom per element was wrong, not atomic theory. Aston received the 1922 Nobel Prize in Chemistry. His central conclusion was charmingly simple: a balance measures a mixture's average weight.
\end{zhengju}

Mass spectrometers have long left the laboratory to become chemistry's busiest balances. Hospitals use a newborn's heel-prick blood drop to screen for dozens of congenital metabolic diseases. Sports testing finds banned drugs at parts-per-billion levels in urine. Airport security wipes your luggage and feeds the test strip into a miniature mass spectrometer to sniff explosive traces. Agricultural testing checks pesticide residues on vegetables. Another medical application directly concerns this chapter: the carbon-13 breath test. A doctor gives you a capsule containing carbon-13-labeled urea---carbon-13 is entirely stable and nonradioactive. If Helicobacter pylori inhabits your stomach, the bacteria break down the urea, producing carbon-13-containing carbon dioxide in your breath that the instrument detects by weight. Using a differently weighted version of the same element as a tracking label is one of isotopes' most elegant applications.

\begin{shiyan}
You can make a tabletop analogy of a mass spectrometer at home. Materials: a hair dryer, a table-tennis ball, a glass marble or solid metal ball, a flat table, and chalk. Procedure: (1) Draw a straight chalk line from one end of the table and let balls roll naturally along it. (2) Secure the dryer beside the table, blowing a steady crosswind perpendicular to the line to make a wind zone. (3) Roll the table-tennis ball and marble through the zone at similar speeds and observe their departures from the line. What to observe: which ball bends more? The light table-tennis ball deflects visibly while the heavy marble nearly follows a straight line, analogous to lighter ions bending more and heavier ions less. Two reminders: this is only an analogy; real ions are accelerated by electric fields and deflected magnetically, without anyone blowing on them. For safety, do not aim the dryer at people or flammable objects, and never handle electrical appliances with wet hands.
\end{shiyan}

\section{Average Weight: Relative Atomic Mass}

With mass spectrometry, decimal atomic masses lose their mystery. Chemists take a carbon-12 atom as the standard and define $\frac{1}{12}$ of its mass as one unit. An atom's mass measured in these units is its relative mass. The table's number, such as chlorine's 35.45, is the \textbf{abundance-weighted average of all the element's natural isotopes}, called its \textbf{relative atomic mass}.

\begin{qiaoliang}
Calculate chlorine's average. Mass spectrometry shows natural chlorine as about 75\% chlorine-35 and 25\% chlorine-37. In a large random collection, roughly three of every four atoms are 35 and one is 37, so the average is about

\[\frac{35\times 3 + 37\times 1}{4} = \frac{105+37}{4} = 35.5\]

This nearly matches the table's 35.45. The remaining difference comes from exact isotope masses being slightly below integer mass numbers; we will not worry about it here. A frequent question deserves emphasis: no chlorine atom in nature has a mass of 35.45, not even one. That is a mixture's average, just as 1.8 children per household does not mean someone actually has 0.8 of a child.
\end{qiaoliang}

Try a real order-of-magnitude estimate of carbon-14's scarcity. A 70-kilogram adult is about 18\% carbon, or 12.6 kilograms. Each \ud{12}{g} contains about $6\times 10^{23}$ atoms (Chapter 5's mole), so the body contains about $6\times 10^{26}$ carbon atoms. Living organisms' carbon-14 fraction is roughly $1.2\times 10^{-12}$, a little over one atom per trillion carbons. Multiplication gives about $7\times 10^{14}$ carbon-14 atoms in you now. Seven hundred trillion sounds alarming, but they are diluted among six hundred thousand septillion ordinary carbon atoms, and only a few thousand spontaneously change per second on average. You live with a whole isotope zoo inside, without feeling a thing.

\section{The Limits of the Identity Card}

Every model has boundaries, including this identity card.

First, proton number locks identity within a clear domain: chemical reactions. Those rearrange external electrons without even knocking on the nuclear door, so lead remains lead throughout burning, boiling, dissolving, and electrifying. At energies high enough to alter nuclei themselves---particle bombardment, fusion, or fission---proton counts can change, and lead really can become gold. Scientists have done it in accelerators, at electricity and equipment costs vastly exceeding the gold's value. Alchemists failed because their furnaces could not reach nuclei, not because identity rules were wrong. Chapter 8's whole story lies beyond this boundary.

Second, an identity card says who, not what its owner will do. Pure carbon arranged in sheets makes soft, slippery graphite; in a three-dimensional network, hard diamond; in hollow balls, fullerenes. Identical elemental identity, entirely different substances. Eight-proton oxygen atoms paired as \chem{O_2} sustain breathing; grouped in threes as \chem{O_3}, they make pungent ozone. To know a substance's personality, identity alone is insufficient. We need its atomic connections, arrangements, and electron exchanges, the subject of the next part on chemical bonds.

Third, never drop ``nearly'' from isotopes' nearly identical chemistry. Most elements differ negligibly, but hydrogen is so light that deuterium's doubled mass makes textbook-worthy differences. Heavy water participates in most ordinary-water reactions but significantly more slowly, and organisms drinking only heavy water over time develop serious problems. The lighter the element, the less its isotope differences can be ignored.

\begin{wuqu}
``Isotopes are dangerous radioactive atoms.'' Frequent news references to radioactive isotopes have encouraged people to equate the two. In fact, isotopes simply share proton number and differ in neutron number; radioactivity does not necessarily follow. Your water contains oxygen-16, oxygen-17, and oxygen-18. Your body's carbon and potassium come with natural isotope mixtures, most stable for billions of years. Radioactive isotopes that decay spontaneously are a minority, and danger also depends on dose and exposure route (Chapter 8 explains). Even pregnant women and children can undergo carbon-13 breath testing because it uses a stable isotope. Before worrying at the word isotope, ask: is it stable, and what is the dose?
\end{wuqu}

\section{Back to the Pencil Point}

We can now keep our two opening promises.

Why is it a lead pencil? In the sixteenth century, the English found a graphite deposit, saw that its shiny gray-black material marked paper, and mistook it for lead ore. The mistaken name persisted. But check the identity card: graphite nuclei have 6 protons, carbon; lead nuclei have 82, lead. Names can deceive; proton counts cannot. Modern reverse deceptions exist too: gold-plated jewelry and diamond-imitation zircon may fool the eye, but atomic numbers reveal them.

Why cannot lead become gold? Three protons separate them, and fires, reagents, and electric currents rearrange electrons without reaching nuclei. The real value of centuries of alchemical failure was forcing humanity to recognize that chemical reactions alter combinations, not identities. Rewriting identity requires knocking on the nuclear door.

Check your own body too. Iron ions in your blood and an iron nail are the same element. Potassium ions in cells and potassium metal that explodes on meeting water are the same element. Calcium ions in bones and calcium in lime are the same element. You are a collection of elements carrying identity cards and different charge states, assembled through precise connections. The next eighty-odd chapters read that assembly manual page by page.

\section{Someone Else Is Waiting at the Door}

Several times this chapter has stopped halfway through carbon-14's story. Its two extra neutrons prevent permanent stability. On average, after a little over fifty-seven hundred years, half a batch of carbon-14 atoms has transformed. This half is a matter of probability, not convention. Nobody can predict which atom will change in which second, yet trillions together behave with clocklike precision. This combination of unpredictable individuals and punctual populations lets archaeologists date an animal bone, doctors attack tumors with radiation, and humans release immense energy sleeping in nuclei. Next chapter opens the nuclear door we have kept pretending was not there.

\begin{tiaozhan}
Skip this passage without losing the main story. Attentive readers may spot a logical trap: relative atomic mass uses $\frac{1}{12}$ of carbon-12's mass as its standard, while mass spectrometers measure isotope masses. What calibrates the spectrometer? Scientific units need an anchor. Before 2019, the SI kilogram was defined by a platinum--iridium cylinder outside Paris, while the mole was tied to the atoms in \ud{12}{g} of carbon-12. In 2019, the international conference on weights and measures made a major change, replacing physical-object anchors with natural constants. It fixed the Avogadro constant directly at $N_A = \ud{6.02214076\times 10^{23}}{mol^{-1}}$, while the kilogram was derived from the Planck constant. Carbon-12's molar mass ceased to be exactly \ud{12}{g/mol} and became a measured value with tiny uncertainty. Your measured weight did not change by an ounce, but science's anchor moved from a wearing metal lump to constants of the universe. A small problem remains: boron's relative atomic mass is 10.81, and natural boron contains only boron-10 and boron-11. Can you infer their percentages? (Let boron-11's fraction be $x$. Solve $11x+10(1-x)=10.81$ to obtain $x=0.81$, about 81\%.)
\end{tiaozhan}

\section*{Try It Yourself}

\begin{sikaoti}
\item Draw three boxes for nuclei containing 6 protons and 6 neutrons, 6 protons and 8 neutrons, and 8 protons and 8 neutrons. Draw protons as dots marked $+$ and neutrons as empty circles. Label each element and isotope using ${}^{A}_{Z}\chem{X}$. Which two atoms have nearly identical chemistry, and why?
\item An oxygen atom with 8 protons, 8 neutrons, and 8 electrons gains two electrons to become \chem{O^{2-}}. List each particle count before and after, and explain why it remains oxygen.
\item Natural copper contains only copper-63 and copper-65, and its relative atomic mass is 63.55. Without looking anything up, decide which isotope is more abundant and explain. Then let copper-63's fraction be $x$, write an equation, and calculate the percentages.
\item Explain mass spectrometry to a friend using the hair-dryer, table-tennis-ball, and marble analogy. Why must ions first carry charge, and why do lighter ions bend more? Where does the analogy fail? (Hint: magnetic fields do not act on stationary charges, so ions must first be accelerated.)
\item An advertisement claims a product turns atmospheric oxygen into healthier negative oxygen ions. Apply this chapter's identity logic: does \chem{O_2} gaining electrons change its elemental identity? Is the resulting substance the same as before? Which two concepts does the advertisement confuse?
\end{sikaoti}
